\section{Appendix C: A worked retrospective on Compound Proposal 289}
\label{sec:appendix-c}

This appendix walks through what the Chio substrate would have done if the Compound DAO had been structured as a Chio polity at the time of Proposal~289 (July 2024). The case study is included to answer two questions a hostile reader will ask: would a Chio polity have \emph{prevented} the attack, and would it have made the recovery cheaper or faster? The honest answers are no and yes, in that order.

\paragraph{The incident.}
On 24 July 2024 the wallet cluster colloquially known as ``Golden Boys'' / ``Humpy'' filed Compound governance Proposal~289 to move approximately 499{,}000 COMP (\$24M at then-prevailing prices) from the Compound DAO treasury to a contract under its control \cite{compoundProposal289}. The proposal passed the 28-29 July vote by 682{,}191 to 633{,}636 (52\% margin) under low voter turnout, with 57 wallets participating. A counter-proposal and negotiated settlement on 30 July rescinded the transfer.

\paragraph{Q1: Which Chio predicates would have failed admission?}
A Chio-structured Compound constitution would have carried at minimum four predicates in $K$: (a) a \emph{treasury-cap} predicate bounding any single-proposal transfer below a constitutional threshold; (b) a \emph{sibling-sum} predicate over delegate concentration, denying when one or two wallet clusters together exceed a constitutional fraction of voting weight; (c) a \emph{participation-floor} predicate denying when fewer than a constitutional minimum of distinct citizen-DIDs participate; (d) a \emph{cross-attempt deduplication} predicate denying repeat-by-cluster within a cooling-off window. Proposal~289 failed (b) and (c) on the actual vote: 57 wallets is below any reasonable participation floor, and the Golden Boys cluster's effective delegation share exceeded the safe sibling-sum bound. The attack would have been admitted under condition (a) alone (the transfer was technically a vote-approved amendment), but the joint predicate fails on (b) and (c).

\paragraph{Q2: Was Proposal 289 structured as a refinement?}
Mostly no. Treasury transfers in Compound governance are conventionally amendments by precedent rather than refinement-checked constitutional changes; they widen the discretion-set of the governance contract for the duration of the transfer. Under Chio's backward-refinement rule, an amendment is admissible only when $\forall r.\ K'(r) \Rightarrow K(r)$ on receipts already admitted under $K$. A treasury-withdrawal proposal that leaves the treasury empty does not narrow $K$ for prior-admitted receipts; it widens executive discretion. Under Chio's enactment invariant (\thm{amendment_admissible_iff_backward_refinement} and the type-level requirement that \emph{enactAmendment} consume a \emph{ConstitutionalDelta} carrying a \emph{BackwardRefines} proof), the vote could have passed politically but enactment would have been hard-rejected.

\paragraph{Q3: Would refinement framing have admitted the attack via a constitutional ratchet?}
The harder case. A sophisticated attacker frames Proposal~289 as ``narrow the set of treasury-eligible recipient addresses to $\{A, B, C\}$,'' which is technically a refinement (the discretion-set shrinks). A constitutional-ratchet attack composes individually valid refinements into an adversary-only predicate; the essential-predicate invariant patches this by requiring that the constitution carry an \emph{essential} predicate set whose admitted-receipt set is preserved across all amendment trajectories. With an essential-predicate invariant including (b)~sibling-sum and (c)~participation-floor, the narrowing-refinement attack vector still fails at enactment.

\paragraph{Q4: What evidence would Chio have produced that current on-chain governance does not?}
Chio's substrate produces, for each rejected proposal: a signed denial receipt naming the predicate(s) violated, the ladder-floor in force, the citizenship-roster snapshot, and the verifier-store identity. Gauntlet's risk analysis and the bribery-wallet flow analysis that catalysed the rescission vote would still have to happen, but they would have a tamper-evident audit corpus to anchor on -- and the rescission itself, as a counter-amendment, would carry its own refinement proof.

\paragraph{Q5: Would Chio have prevented the attack?}
No. The substrate makes the attack \emph{legible at admission} rather than preventing it. A Compound DAO that adopted Chio without also adopting the essential-predicate invariant and a participation-floor predicate would have admitted Proposal~289 just as the historical vote did. The substrate's contribution is that admission is auditable in a canonical format, the rescission has a refinement proof, and the bribery wallet flows produce signed denial receipts rather than off-chain forensic narratives. This is a real improvement, but it is not prevention.

\paragraph{Implications for Chio adoption.}
Compound 289 exemplifies why \S\ref{sec:limits}'s ``Amendment expressiveness'' bullet is load-bearing: backward refinement preserves admitted history, but real governance failures are usually trajectory failures across many individually-defensible refinements. The substrate's enforceable contribution is the artifact path, not the political one; the political work (essential predicates, participation thresholds, sibling-sum bounds) is still on the DAO to specify.
